% Part:sets-functions-relations
% Chapter: size-of-sets
% Section: non-enumerability

\documentclass[../../../include/open-logic-section]{subfiles}

\begin{document}

\olfileid{sfr}{siz}{nen}

\olsection{\printtoken{S}{nonenumerable} संच}

\begin{editorial}
  या विभागात \olref[enm]{sec} मधील व्याख्या वापरून $\Bin^\omega$ आणि
  $\Pow{\PosInt}$ यांची अगणनीयता सिद्ध केली आहे. ही मांडणी
  \olref[enm-alt]{sec} मधील आवृत्तीपेक्षा थोडी अधिक प्राथमिक आणि
  थोडी अधिक तपशीलवार ठेवली आहे.
\end{editorial}

धन पूर्णांकांचा संच~$\PosInt$ यासारखे काही संच अनंत असतात. आतापर्यंत
पाहिलेल्या अनंत संचांची सर्व उदाहरणे !!{enumerable} होती. पण हा गुणधर्म
नसलेले अनंत संचही असतात. अशा संचांना \emph{!!{nonenumerable}} म्हणतात.

सुरुवातीलाच, !!{nonenumerable} संच असतात ही गोष्ट कदाचित आश्चर्यकारक
वाटेल. कोणत्याही !!{enumerable} संच~$A$ साठी $f \colon \PosInt \to A$
असे !!a{surjective} फलन असते. संच !!{nonenumerable} असेल, तर असे फलन
असू शकत नाही. म्हणजे, $\PosInt$ मधील अनंत !!{element}s ना~$A$ कडे
पाठवणारे कोणतेही फलन~$A$ मधील सर्व घटक व्यापू शकत नाही. या अर्थाने,
अनंत असलेल्या धन पूर्णांकांपेक्षा~$A$ मध्ये ``अधिक'' !!{element}s असतात.

एखादा संच !!{nonenumerable} आहे हे कसे सिद्ध करायचे? तसे कोणतेही
आच्छादक फलन अस्तित्वात असू शकत नाही, हे दाखवावे लागते. समतुल्यपणे,
$A$ मधील घटकांचे एकदिश अनंत यादीत प्रगणन करता येत नाही, हे दाखवावे
लागते. हे करण्याचा सर्वोत्तम मार्ग म्हणजे~$A$ मधील !!{element}s च्या
प्रत्येक यादीतून किमान एक घटक सुटतो, किंवा $f\colon \PosInt \to A$
असे कोणतेही फलन !!{surjective} असू शकत नाही, हे दाखवणे. यासाठी कँटर
यांची \emph{कर्णरेषीय पद्धत} वापरता येते. $A$ मधील !!{element}s ची
एखादी यादी, म्हणजे $x_1$, $x_2$, \dots, दिली असता, आपण~$A$ मधील असा
दुसरा घटक रचतो की त्याच्या रचनेमुळे तो त्या यादीत असणे अशक्य ठरते.

आपले पहिले उदाहरण म्हणजे $0$ आणि $1$ यांच्या सर्व अनंत व मध्ये खंड
नसलेल्या अनुक्रमांचा संच~$\Bin^\omega$.

\begin{thm}
\ollabel{thm:nonenum-bin-omega}
$\Bin^\omega$ हा !!{nonenumerable} आहे.
\end{thm}

\begin{proof}
व्याघाताने सिद्ध करण्यासाठी असे समजा की $\Bin^\omega$ हा !!{enumerable} आहे; म्हणजे,
$\Bin^\omega$ मधील सर्व !!{element}s ची $s_{1}$, $s_{2}$, $s_{3}$,
$s_{4}$, \dots{} अशी यादी आहे, असे समजा. यांपैकी प्रत्येक $s_i$ हा
स्वतः $0$ आणि~$1$ यांचा अनंत अनुक्रम आहे. या यादीतील $i$ व्या अनुक्रमाच्या
$j$ व्या घटकाला $s_i(j)$ म्हणू. मग $i$ वा अनुक्रम~$s_i$ असा आहे:
\[
s_i(1), s_i(2), s_i(3), \dots
\]

ही यादी आणि तिच्यातील प्रत्येक अनुक्रम~$s_i$ चे घटक आपण पुढील सरणीत
मांडू शकतो:
\[
\begin{array}{c|c|c|c|c|c}
& 1 & 2 & 3 & 4 & \dots \\\hline
1 & \mathbf{s_{1}(1)} & s_{1}(2) & s_{1}(3) & s_1(4) & \dots \\\hline
2 & s_{2}(1)& \mathbf{s_{2}(2)} & s_2(3) & s_2(4) & \dots \\\hline
3 & s_{3}(1)& s_{3}(2) & \mathbf{s_3(3)} & s_3(4) & \dots \\\hline
4 & s_{4}(1)& s_{4}(2) & s_4(3) & \mathbf{s_4(4)} & \dots \\\hline
\vdots & \vdots & \vdots & \vdots & \vdots & \mathbf{\ddots}
\end{array}
\]
डाव्या बाजूकडील खुणा यादीतील अनुक्रमांना $s_1$, $s_2$, \dots{} असे
क्रमांक देतात; वरच्या बाजूकडील संख्या प्रत्येक अनुक्रमातील !!{element}s
ना क्रमांक देतात. उदाहरणार्थ, $s_{1}(1)$ हे अनुक्रम~$s_{1}$ मधील
पहिला !!{element} असलेल्या संख्येचे नाव आहे; ती संख्या $0$ किंवा~$1$
असू शकते. पुढेही याचप्रमाणे अर्थ घ्यायचा.

आता आपण $0$ आणि $1$ यांचा असा अनंत अनुक्रम~$\overline{s}$ रचू, जो
या यादीत असणे अशक्य आहे. $\overline{s}$ ची व्याख्या $s_1$, $s_2$,
\dots{} या यादीवर अवलंबून असेल. $0$ आणि $1$ यांच्या अनंत अनुक्रमांची
कोणतीही अनंत यादी असा अनंत अनुक्रम~$\overline{s}$ निर्माण करते, जो
त्या यादीत नसेल याची खात्री असते.

$\overline{s}$ ची व्याख्या करण्यासाठी आपण त्याचे सर्व !!{element}s
ठरवतो; म्हणजे, प्रत्येक $n \in \PosInt$ साठी $\overline{s}(n)$
ठरवतो. वरील सरणीच्या कर्णरेषेवरून खाली वाचून (म्हणूनच ``कर्णरेषीय
पद्धत'' हे नाव) प्रत्येक $1$ चे $0$ मध्ये आणि प्रत्येक $0$ चे~$1$
मध्ये रूपांतर करून आपण हे करतो. अधिक अमूर्तपणे, कर्णरेषेवरील $n$ वा
!!{element}, म्हणजे $s_n(n)$, हा $1$ आहे की $0$ यानुसार
$\overline{s}(n)$ हे $0$ किंवा $1$ असे परिभाषित करतो:
\[
\overline{s}(n) =
\begin{cases}
1 & \text{जर $s_{n}(n) = 0$ असेल}\\
0 & \text{जर $s_{n}(n) = 1$ असेल}.
\end{cases}
\]
प्रकरणांनुसार व्याख्येपेक्षा सूत्रे अधिक आवडत असतील, तर
$\overline{s}(n) = 1 - s_n(n)$ अशीही व्याख्या करता येईल.

$\overline{s}$ हा स्पष्टपणे $0$ आणि $1$ यांचा अनंत अनुक्रम आहे, कारण
तो आपल्या सरणीच्या कर्णरेषेवर दिसणाऱ्या $0$ आणि $1$ यांच्या अनुक्रमाचा
पूरक अनुक्रम आहे. म्हणून $\overline{s}$ हा $\Bin^\omega$ चा
!!a{element} आहे. पण तो $s_1$, $s_2$, \dots{} या यादीत असू शकत नाही.
असे का?

तो यादीतील पहिला अनुक्रम~$s_1$ असू शकत नाही, कारण त्याचा पहिला
!!{element} हा $s_1$ च्या पहिल्या घटकापेक्षा वेगळा आहे. $s_1(1)$
काहीही असो, $\overline{s}(1)$ त्याच्या विरुद्ध असेल अशी आपण व्याख्या
केली. तो यादीतील दुसरा अनुक्रमही असू शकत नाही, कारण $\overline{s}$
चा दुसरा घटक $s_2$ च्या दुसऱ्या घटकापेक्षा वेगळा आहे: $s_2(2)$ हा
$0$ असेल, तर $\overline{s}(2)$ हा $1$ आहे, आणि उलट. पुढेही असेच.

अधिक नेमकेपणाने सांगायचे तर, $\overline{s}$ यादीत असता, तर कोणत्या
तरी $k$ साठी $\overline{s} = s_{k}$ असते. दोन अनुक्रम प्रत्येक स्थानी
जुळतात तेव्हा आणि केवळ तेव्हाच ते एकरूप असतात; म्हणजे, कोणत्याही~$n$
साठी $\overline{s}(n) = s_{k}(n)$ असते. म्हणून, विशेष प्रकरण म्हणून
$n = k$ घेतल्यास $\overline{s}(k) = s_{k}(k)$ असणे आवश्यक ठरते.
$s_k(k)$ हे $0$ किंवा~$1$ आहे. ते $0$ असेल, तर $\overline{s}(k)$ हे~$1$
असले पाहिजे---कारण आपण $\overline{s}$ ची तशीच व्याख्या केली आहे. पण
$s_k(k) = 1$ असेल, तर पुन्हा $\overline{s}$ च्या व्याख्येमुळे
$\overline{s}(k) = 0$ होते. दोन्ही बाबतीत
$\overline{s}(k) \neq s_{k}(k)$.

सुरुवातीला आपण $\Bin^\omega$ मधील !!{element}s ची $s_1$, $s_2$,
\dots{} अशी यादी आहे असे गृहीत धरले. या यादीपासून आपण असा
अनुक्रम~$\overline{s}$ रचला की तो यादीत असू शकत नाही, हे सिद्ध केले.
पण सर्व $s_i$ हे $0$ आणि $1$ यांचे अनुक्रम असतील, तर तो नक्कीच $0$
आणि $1$ यांचा अनुक्रम असतो; म्हणजे, $\overline{s} \in
\Bin^\omega$. यावरून विशेषतः $\Bin^\omega$ मधील \emph{सर्व}
!!{element}s ची यादी असू शकत नाही. कारण अशी कोणतीही यादी दिली, तरी
तिच्यात नसेल असा अनुक्रम~$\overline{s}$ आपण रचू शकतो. म्हणून
$\Bin^\omega$ मधील सर्व अनुक्रमांची यादी आहे हे गृहीतक व्याघात निर्माण
करते.
\end{proof}

\begin{explain}
या सिद्धता-पद्धतीला ``कर्णीकरण'' म्हणतात, कारण तिच्यात सरणीच्या
कर्णरेषेचा वापर करून~$\overline{s}$ ची व्याख्या केली जाते. कर्णीकरणात
सरणी असलीच पाहिजे असे नाही: प्रत्यक्ष सरणी किंवा कर्णरेषा नसतानाही
हीच कल्पना वापरून संच !!{enumerable} नाहीत, हे दाखवता येते.
\end{explain}

\begin{thm}
\ollabel{thm:nonenum-pownat}
$\Pow{\PosInt}$ हा !!{enumerable} नाही.
\end{thm}

\begin{proof}
आपण त्याच पद्धतीने पुढे जाऊ: $\PosInt$ च्या उपसंचांची प्रत्येक यादी
दिली असता त्या यादीत नसलेला $\PosInt$ चा एक उपसंच असतो, हे दाखवू.
$\PosInt$ च्या उपसंचांची पुढील यादी दिली आहे असे समजा:
\[
Z_{1}, Z_{2}, Z_{3}, \dots
\]
आता आपण $\overline{Z}$ हा असा संच परिभाषित करतो की प्रत्येक
$n \in \PosInt$ साठी $n \in \overline{Z}$ तेव्हा आणि केवळ तेव्हाच,
जेव्हा $n \notin Z_{n}$:
\[
\overline{Z} = \Setabs{n \in \PosInt}{n \notin Z_n}
\]
$\overline{Z}$ हा स्पष्टपणे धन पूर्णांकांचा संच आहे, कारण गृहीतकानुसार
प्रत्येक~$Z_n$ तसा आहे; म्हणून $\overline{Z} \in
\Pow{\PosInt}$. पण $\overline{Z}$ यादीत असू शकत नाही. हे दाखवण्यासाठी
प्रत्येक $k \in \PosInt$ साठी $\overline{Z} \neq Z_k$, हे सिद्ध करू.

आता $k \in \PosInt$ स्वैर असू द्या. प्रत्येक $n \in \PosInt$ साठी
$n \in \overline{Z}$ तेव्हा आणि केवळ तेव्हाच, जेव्हा $n \notin Z_n$,
अशी $\overline{Z}$ ची व्याख्या केली आहे. विशेषतः $n=k$ घेतल्यास,
$k \in \overline{Z}$ तेव्हा आणि केवळ तेव्हाच, जेव्हा $k \notin Z_k$.
पण यावरून $\overline{Z} \neq Z_k$ हे दिसते, कारण $k$ हा एकाचा
!!a{element} आहे आणि दुसऱ्याचा नाही; म्हणून $\overline{Z}$ आणि $Z_k$
यांचे !!{element}s वेगवेगळे आहेत. $k$ स्वैर असल्याने $\overline{Z}$ हा
$Z_1$, $Z_2$, \dots{} या यादीत नाही.
\end{proof}

\begin{explain}
मागील सिद्धतेत कर्णरेषेचा उल्लेख झाला नाही; पण पुढील चित्र मनात आणले,
तर तिच्यात कर्णरेषा आहे असे समजता येते. $Z_1$, $Z_2$, \dots{} हे संच
एका सरणीत लिहिले आहेत अशी कल्पना करा; प्रत्येक !!{element}~$j \in Z_i$
हा $j$ व्या स्तंभात लिहिला आहे. समजा, त्या यादीतील पहिले चार संच
$\{1,2,3,\dots\}$, $\{2, 4, 6, \dots\}$, $\{1,2,5\}$ आणि
$\{3,4,5,\dots\}$ आहेत. मग सरणीची सुरुवात अशी होईल:
\[
\begin{array}{r@{}rrrrrrr}
  Z_1 = \{ & \mathbf{1}, & 2, & 3, & 4, & 5, & 6, & \dots\}\\
  Z_2 = \{ &  & \mathbf{2}, &  & 4, &  & 6, & \dots\}\\
  Z_3 = \{ & 1, & 2, &  &  & 5\phantom{,} &  & \}\\
  Z_4 = \{ &  &  & 3, & \mathbf{4}, & 5, & 6, & \dots\}\\
  \vdots & & & & & \ddots
\end{array}
\]
यानंतर कर्णरेषेवरून खाली जाताना कर्णरेषेवर दिसणाऱ्या संख्या वगळून,
आणि $j$ व्या ओळीच्या त्याच क्रमांकाच्या स्तंभातील जागा रिकामी असेल अशा~$j$
चा समावेश करून, $\overline{Z}$ हा संच मिळतो. वरील उदाहरणात
आपण $1$ आणि $2$ वगळू, $3$ चा समावेश करू, $4$ वगळू, इत्यादी.
\end{explain}

\begin{prob}
कर्णरेषीय युक्तिवादाने $\Pow{\Nat}$ हा !!{nonenumerable} आहे, हे दाखवा.
\end{prob}

\begin{prob}\label{sfr:siz:nen:prob:f-posint}
$f \colon \PosInt \to \PosInt$ अशा फलनांचा संच !!{nonenumerable}
आहे, हे स्पष्ट कर्णरेषीय युक्तिवादाने दाखवा. म्हणजे, $f_1$, $f_2$,
\dots{} ही फलनांची यादी असेल आणि प्रत्येक $f_i\colon
\PosInt \to \PosInt$ असेल, तर या यादीत नसलेले कोणते तरी
$\overline{f}\colon \PosInt \to \PosInt$ असते, हे दाखवा.
\end{prob}

\end{document}
